%% transmon-benchmark.17 -- revision of v16, 2026-10-16, against the Opus
%% (41/44), Fable (40/44) and Codex reviews of v16, with the grounded-box
%% capacitance bound and the adaptive run on the final design. Text, table and
%% bibliography changes only; no solver was run for this revision. See
%% changelog.md, provenance.md and validation.md. v1-v16 and their critic
%% siblings are unchanged.

\documentclass{anvil-paper}
\usepackage{xurl}
\usepackage{array}
\newcommand{\Cisl}{C_{\mathrm{isl}}}

%% --- URL line-breaking: allow breaks at hyphens and add stretch so the
%% bibliography stays inside the measure (render-gate overfull check).
\expandafter\def\expandafter\UrlBreaks\expandafter{\UrlBreaks\do\-}
\Urlmuskip=0mu plus 1mu
%% Float placement: let text share pages with floats.
\renewcommand{\topfraction}{0.92}
\renewcommand{\bottomfraction}{0.6}
\renewcommand{\textfraction}{0.07}
\renewcommand{\floatpagefraction}{0.85}
\setcounter{topnumber}{3}
\setcounter{totalnumber}{4}
%% Compact bibliography.
\renewcommand{\bibfont}{\fontsize{8}{9}\selectfont}
\setlength{\bibsep}{0pt plus 1pt}

%% --- Draft-phase helper macros -------------------------------------------
%% Figure placeholder: uses the rendered figure if pub-figures has produced
%% it, otherwise a framed placeholder naming the source script.
\newcommand{\anvilfig}[3][0.98]{%
  \IfFileExists{figures/#2}%
    {\includegraphics[width=#1\linewidth]{figures/#2}}%
    {\fbox{\parbox{#1\linewidth}{\centering\vspace{5em}%
      \texttt{figures/#2}\\[0.5em]%
      {\small placeholder: render via \texttt{figures/src/#3}}%
      \vspace{5em}}}}}

\title{Uniform mesh refinement underestimates the frequencies of the
DeviceLayout.jl transmon: an independent finite-element cross-check and
re-optimization}
\author{Robb Walters\\
  {\normalsize Independent Researcher}}
% NOTE: Crutcher Dunnavant to be offered coauthorship at review (operator, 2026-10-13)
\date{October 2026}

\begin{document}
\maketitle

\begin{abstract}
The frequencies that the SingleTransmon example of DeviceLayout.jl
publishes for an initial and an optimized design are
discretization-limited. We reproduce them with Palace within $0.7\%$ and
cross-check Palace with geode-fem, an independent solver;
on shared meshes the two agree to $2\times10^{-4}\%$ on the qubit and
$0.03\%$ on the resonator. All frequencies are linear eigenmodes. Adaptive
Palace runs, still rising, put the qubits $7.0\%$ and $8.0\%$ above their
published values without extrapolation. On the initial design, a two-sided
island-capacitance bound for the grounded box, read through
a circuit model, caps the qubit at $4.477$\,GHz ($+8.1\%$). Extrapolated at order 1 in
mesh size, the fastest rate the sheet edges allow on uniform
meshes, two refinement families put the qubits $7.3$--$9.6\%$ and the
resonators $2.7$--$3.3\%$ higher; these estimates
lean low and rest on Palace's levels alone, since geode-fem ran only order
1 on the two coarsest meshes. At these frequencies the $3.0$/$4.0$\,GHz design
needs a capacitor $19$--$20\%$ longer than published ($17$--$39\%$ within
the band). Counting every Palace run, probes included, Newton's method meets DeviceLayout's
criterion in 3 runs, with geode-fem's shape derivative or a
geode-fem-calibrated scaling law in its Jacobian, and in 4 without
geode-fem; PRIMA takes 10 to 12 in our environment.
\end{abstract}

\section{Introduction}
\label{sec:intro}

Superconducting-qubit layouts are designed by iterating field solves
against frequency targets, a loop that QDesignOptimizer
\citep{eriksson2025automated} steers with analytic models. DeviceLayout.jl \citep{devicelayout} builds the layout and
its mesh, and Palace \citep{palace}, an MFEM-based solver \citep{mfem},
computes the eigenmodes. DeviceLayout's SingleTransmon example
\citep{devicelayout-blog,devicelayout-example} runs the loop end to end. It
publishes an initial transmon-plus-resonator design at $4.14$ and
$5.591$\,GHz and a design optimized for targets of $3.0$ and $4.0$\,GHz, at
$3.005$ and $3.982$\,GHz, found with a derivative-free optimizer in nine
Palace runs. Like every frequency in this paper, these are eigenfrequencies
of a linear model with the junction as a lumped inductor: the qubit
frequency is the junction's linearized (plasma) mode, not the transmon's
0--1 transition $f_{01}$, which lies below it by about the charging energy
$E_C/h$ (Section~\ref{sec:junction}).

Those frequencies come from one discretization, and the example's authors
say what that implies. They run Palace at second order ``to reduce
sensitivity to the discretization'' and note that this ``does not mean that
the eigenfrequencies found by Palace will have converged to their true
values'' \citep{devicelayout-example}. This paper measures
how far from converged the published frequencies are. Our
main result is that the published frequencies are discretization-limited.
Palace's adaptive refinement reaches qubits $7.0\%$ and $8.0\%$ above the
two published values without extrapolation, and its uniform ladders, in two
refinement families, extrapolate to qubits $7.3$--$9.6\%$ above them;
every uniform level we could afford lies below those estimates. geode-fem
reproduces every level of those ladders it can reach, the two coarsest
meshes at order 1, to $10^{-4}\%$ on the qubit, and a
two-sided bound on the island capacitance, read through a circuit model, is
consistent with that shift and caps the initial qubit under that model.

The second code is GEODE-FEM (``geode-fem''), a Rust finite-element solver
on the Burn tensor framework \citep{burn}. It solves the same
$H(\mathrm{curl})$ eigenproblem as Palace with a separate implementation,
and it computes design derivatives, among them the shape derivative of an
eigenfrequency. Given the same input mesh, it reproduces Palace's modes
(Section~\ref{sec:agreement}), and on the $h/2$ meshes of
Section~\ref{sec:crosscheck} the two codes solve the same discrete problem
apart from the junction term. This gives an independent
same-mesh cross-check of the Palace/DeviceLayout.jl transmon stack,
complementing the earlier checks of Palace by SQDMetal
\citep{sommers2025open} and \citet{ye2025electromagnetic}.

Our contributions are:
\begin{itemize}
  \item \textbf{A refinement study of a published design example}, which
    puts its qubits $7.3$--$9.6\%$ above their published values, with a
    second refinement family and adaptive runs on both designs checking
    the qubit estimates; the finer levels and the extrapolations are
    Palace's (Section~\ref{sec:crosscheck}).
  \item \textbf{Two-sided capacitance bounds with the sapphire tensor.}
    An equilibrated-flux lower bound and conforming upper bounds bracket
    the island capacitance $\Cisl$ of the unmodified device, without the
    junction capacitance, in $[79.56, 79.77]$\,fF with the vacuum box
    grounded, as in the eigenproblem, and in $[79.56, 79.76]$\,fF with it
    natural; under a two-node circuit model the lower bound caps the
    initial qubit at $4.477$\,GHz (Section~\ref{sec:static}).
  \item \textbf{A re-optimized design, with controls.} At the extrapolated
    frequencies the $3.0$/$4.0$\,GHz design needs $c \approx
    1451$--$1466\,\mu$m against the published $1217.678\,\mu$m. Newton's
    method, with and without geode-fem input, Broyden's method and PRIMA
    are compared by the Palace runs they need in one environment
    (Section~\ref{sec:reopt}).
  \item \textbf{Same-input-mesh agreement with Palace}, with the
    $17.49$\,GHz mode that both codes share traced to the isotropic sheet
    term of the port model (Section~\ref{sec:agreement}).
\end{itemize}

\section{Related Work}
\label{sec:related}

\paragraph{Cross-checks of Palace.}
SQDMetal \citep{sommers2025open} cross-checks a Palace workflow against
COMSOL and Ansys with mesh-convergence studies, and
\citet{ye2025electromagnetic} compare a Palace workflow with cryogenic
measurement, matching resonator frequencies to within $0.3\%$ with a mesh
refined at the CPW gaps and fourth-order elements, a setting their
convergence study puts within $\pm0.3\%$ of extrapolated values;
energy-participation analyses likewise check convergence over adaptive
passes \citep{minev2021energy}. Both cross-checks use models built
independently for each tool. Our comparison instead holds the input mesh
and the junction model fixed and changes the solver implementation, and it
adds a refinement study of a published design: code-to-code comparison and
solution verification \citep{oberkampf2010verification}, not validation
against measurement. As of October 2026, neither Palace nor DeviceLayout.jl has a
methods publication; we cite their repositories and documentation.

\paragraph{Design sensitivities and qubit-design optimization.}
Adjoint sensitivity analysis of frequency-domain electromagnetic solvers is
an established field \citep{nikolova2004adjoint}, including shape adjoints
for Maxwell cavity eigenvalues \citep{akcelik2005adjoint}; Palace exposes no
derivatives.
In qubit design, SQcircuit \citep{rajabzadeh2023analysis} and
SuperGrad \citep{wang2025supergrad} differentiate circuit Hamiltonians,
SQuADDS \citep{shanto2024squadds} looks designs up in what its authors
call a ``validated design database'', and \citet{eun2023shape} shape-optimize transmon pads with FEM
and a derivative-free optimizer.
DeviceLayout's example optimizes with PRIMA \citep{zhang2023prima}, a
reference implementation of Powell's model-based derivative-free methods;
with bounds only it runs BOBYQA \citep{powell2009bobyqa}.

\paragraph{Two-sided capacitance bounds and circuit models.}
Bracketing a capacitance between complementary energy principles is
classical \citep{specogna2011complementary}. Our lower bound uses the
hypercircle idea of \citet{prager1947approximations} with Raviart--Thomas
fluxes \citep{raviart1977mixed} equilibrated patch by patch
\citep{braess2007equilibrated,ern2015polynomial}. The method is not new; we
apply it to an island capacitance with a floating conductor on a
production qubit layout. A web search in October 2026 of the
complementary-bound and qubit-design literature (queries such as
``guaranteed capacitance bounds qubit'') found no two-sided capacitance
bound in the qubit-design literature. The two-node circuit of Section~\ref{sec:junction} is a minimal instance of
the quasi-lumped route \citep{minev2021circuit}, which partitions a device
into lumped or quasi-distributed cells, simulates each and reduces the
composite to coupled subsystems; the energy-participation route
\citep{minev2021energy} instead quantizes the distributed eigenmodes.

\section{Solvers, Device and Models}
\label{sec:setup}

\subsection{geode-fem}
\label{sec:architecture}

geode-fem assembles first-order N\'ed\'elec \citep{nedelec1980} and scalar
P1/P2 element operators as batched $[n_{\mathrm{elem}},k,k]$ tensor
operations on Burn. Its eigensolver is shift-invert Lanczos \citep{arpack}
on a real symmetric pencil, factorized with faer \citep{kazdadi2026faer}. The two codes share the
formulation, the lumped junction port, the materials and the input mesh,
and geode-fem follows a specification of Palace's operators
(Section~\ref{sec:built}); assembly, the eigensolver and the linear algebra
are separate implementations.

The eigenfrequency derivative used in Section~\ref{sec:reopt} needs no
adjoint solve, because the pencil is symmetric: for a simple eigenpair
$(\lambda, x)$ with $x^{\mathsf T}Mx = 1$,
\begin{equation}
  \frac{d\lambda}{d\theta} \;=\;
  x^{\mathsf T}\Bigl(\frac{\partial K}{\partial\theta}
  - \lambda\,\frac{\partial M}{\partial\theta}\Bigr)x ,
  \label{eq:hf-eig}
\end{equation}
the classical eigenvalue derivative of a symmetric generalized eigenproblem
\citep{foxkapoor1968rates}. Here $K$ and $M$ include the junction terms of
Eq.~\eqref{eq:port}; the junction does not move under the capacitor-length
motion of Section~\ref{sec:reopt}, so those terms do not contribute.
geode-fem evaluates $\partial K/\partial\theta$ and $\partial
M/\partial\theta$ along a node velocity $\partial X/\partial\theta$ with a
hand-written forward-mode dual-number twin of its N\'ed\'elec kernel
(tested against it to $10^{-10}$); Burn's autodiff tape is not used.

\subsection{Device, mesh and materials}
\label{sec:geometry}

The device is DeviceLayout.jl v1.15.0's SingleTransmon example
\citep{devicelayout}: a transmon \citep{koch2007charge} with a meander readout resonator and a
feedline on sapphire, in a vacuum box (Figure~\ref{fig:geometry}). Its two
design parameters are the transmon's \texttt{capacitor\_length} $c$ and the
resonator's \texttt{total\_length} $L$; the published initial design has
$c = 620\,\mu$m and $L = 5000\,\mu$m. The sapphire permittivity is
$\mathrm{diag}(9.3, 9.3, 11.5)$ with the optic axis normal to the chip;
the configuration rotates the material axes about that normal, which
leaves a uniaxial tensor unchanged, so the lab tensor is diagonal. Where a
scalar model is used it is the trace average $\bar\varepsilon = 10.03$.

\begin{figure}[htbp]
  \centering
  \anvilfig{fig1-geometry.pdf}{fig1\_geometry.py}
  \caption{Benchmark geometry, drawn from the SHA-256-pinned MSH~4.1
  input mesh of the DeviceLayout.jl v1.15.0 \texttt{SingleTransmon} device.
  (a) Top view of the three metal conductors of Section~\ref{sec:static}:
  the feedline (teal), the qubit island (blue), and the ground plane with
  the meander resonator (gray); squares mark the two readout ports. (b) The
  $1\times1\,\mu$m junction \texttt{lumped\_element} sheet with its
  four-triangle triangulation. (c) Material cross-section, schematic and not
  to scale. Both solvers read this file (22{,}684 nodes; 133{,}314
  tetrahedra).}
  \label{fig:geometry}
\end{figure}

Two kinds of mesh appear. The \emph{fixture} is a first-order mesh of the
initial geometry written by an earlier DeviceLayout.jl v1.15.0 driver,
\path{reference/julia/generate_transmon_fixture.jl} (22{,}684 nodes,
133{,}314 tetrahedra), pinned by SHA-256;
Section~\ref{sec:agreement} and the capacitance brackets of
Section~\ref{sec:static} use it and its refinements, which leave its
polygonal surfaces unchanged. For the cross-check and the re-optimization,
DeviceLayout meshes each geometry with the two parameters exposed, at first
order for the ladders and at its own second order for the runs that
reproduce its discretization; curving the faces changes the qubit by at
most $0.0001\%$ and the resonator by about $0.05\%$. On the
initial geometry the fixture and DeviceLayout's mesh give the same Palace
qubit at $h$ ($3.4527$ against $3.4520$\,GHz, $0.02\%$). PEC is applied on
the metal sheets and the box; the readout ports are left open in every
eigenproblem except the DeviceLayout-configuration runs of
Section~\ref{sec:crosscheck}.

\subsection{Junction and circuit models}
\label{sec:junction}

The Josephson junction is a linear lumped element, $L_J = 14.860$\,nH in
parallel with $C_J = 5.5$\,fF, on the four-triangle sheet $\Gamma$ of
Figure~\ref{fig:geometry}b, as in DeviceLayout's Palace configuration. With
port length $\ell$ and width $w$ the uniform lumped-port construction adds
\begin{equation}
  K_{\mathrm{port}} = \frac{\ell}{w\,L_J}\, S_\Gamma
  \qquad\text{and}\qquad
  M_{\mathrm{port}} = \frac{C_J\,\ell}{w}\, S_\Gamma
  \label{eq:port}
\end{equation}
to the pencil $(K + K_{\mathrm{port}})x = \omega^2(M + M_{\mathrm{port}})x$,
where $S_\Gamma$ is the isotropic surface matrix on $\Gamma$. Palace's
\texttt{LumpedPort} adds the same terms with the same scalings. geode-fem's
rank-1 port-voltage shunt replaces $S_\Gamma$ by $b\,b^{\mathsf T}$, with
$b_i = (1/w)\int_\Gamma N_i\cdot\hat e\,dS$ the port-voltage functional; on
the fixture it moves the qubit by $6\times10^{-5}\,\%$, keeps the other
physical modes to four decimals, removes the
artifact mode of Section~\ref{sec:agreement}, and is the junction model of
every geode-fem eigensolve in Sections~\ref{sec:crosscheck}
and~\ref{sec:reopt}.

The electrostatic model splits the metal into three conductors: ground
with the resonator, the floating feedline, and the island. Its island
capacitance with the feedline floating and neutral,
$\Cisl = C_{ii} - C_{if}^2/C_{ff}$, excludes the junction. In the
convention of \citet{koch2007charge} the transmon's total capacitance is
$C_\Sigma = \Cisl + C_J$, and it sets the charging energy $E_C =
e^2/2C_\Sigma$. A two-node circuit connects it to the eigenproblem: a bare
qubit $1/(2\pi\sqrt{(L_J + L_s)C_\Sigma})$, with $L_s$ the series
inductance of the junction's current path, coupled with coefficient $k$
to the resonator. At $L_s = k = 0$ this is the plasma frequency
$\sqrt{8E_JE_C}/h$, and the transmon's $f_{01} \approx \sqrt{8E_JE_C}/h -
E_C/h$ lies below it. On the fixture $L_s$ and $k$ were fitted on three and
two discretizations; extrapolated, they give bands of $0.113$--$0.125$\,nH
and $0.0227$--$0.0310$, which are estimates.

\section{Same-Input-Mesh Agreement With Palace}
\label{sec:agreement}

Both codes read the fixture with the same materials and junction model at
first order (Table~\ref{tab:agreement}). Palace's configuration schema
defines \texttt{Target} as the frequency above which its eigensolver
searches \citep{palace}, so the six modes above $4.5$\,GHz come from a run
with \texttt{Target}~$=4.5$\,GHz and the qubit from one with
\texttt{Target}~$=2$\,GHz. All six agree to $\leq 0.033\%$, and the qubit mode to $5\times10^{-5}\,\%$, with port
participation $0.9942$ in geode-fem and port EPR $0.99424$ in Palace, the
energy-participation ratio of the junction, the fraction of the mode's
inductive energy stored in it \citep{minev2021energy}. Both codes use the
isotropic sheet term of Eq.~\eqref{eq:port} here.

The two codes read the same input mesh but do not solve the same unknowns.
Palace cracks each metal sheet, but not the junction sheet, into two sides
by duplicating its vertices. Where an edge inside a sheet joins two outline
vertices, which are not duplicated, the sides cannot be separated, and
Palace bisects the elements on both sides of that edge, in two passes on
the fixture (``local bisection
for under-resolved interior boundaries'' in its log;
\texttt{palace/utils/geodata.cpp}). Such edges lie where a sheet is one
element wide, almost all on the resonator and feedline sheets
(Table~\ref{tab:agreement}), so the added tetrahedra are concentrated
there, though some touch the island. Cracking alone leaves the PEC problem unchanged: each
volume edge attaches to one copy of a duplicated vertex, so the edges that
cracking adds are copies of edges in a sheet, and the PEC condition
eliminates them. Only the bisection changes the unknowns.

\begin{table}[htbp]
  \centering
  \caption{Eigenmode agreement from the identical input mesh (22{,}684
  nodes, 133{,}314 tetrahedra), first order in both solvers, junction as
  the lumped $L_J \parallel C_J$ shunt. Counted the same way, as all
  first-order N\'ed\'elec unknowns (edges) including those on PEC sheets,
  geode-fem's mesh has $156{,}863$ ($133{,}108$ interior) and Palace's has
  $183{,}891$, after it cracks the metal sheets and adds $4{,}913$ elements
  on them by bisection, triggered by $1{,}007$ edges, $993$ of them on
  the resonator's ground plane and the feedline. Both
  codes are printed to geode-fem's archived three decimals (six for the
  qubit); $\Delta$ is the archived relative difference, computed from
  Palace's full-precision value. Sources:
  \texttt{benchmarks/transmon\_eigen/results.toml} and
  \texttt{benchmarks/transmon\_eigen/palace\_crack/}.}
  \label{tab:agreement}
  \footnotesize
  \begin{tabular}{@{}llccc@{}}
    \toprule
    Mode & & geode-fem (GHz) & Palace (GHz) & $\Delta$ (\%) \\
    \midrule
    qubit & transmon qubit & 3.452751 & 3.452749 & 0.00005 \\
    resonator & readout resonator & 5.153 & 5.151 & 0.032 \\
    mode 2 & box/cavity & 15.465 & 15.461 & 0.029 \\
    port sheet & port-sheet artifact & 17.490 & 17.490 & 0.001 \\
    mode 4 & box/cavity & 18.693 & 18.692 & 0.007 \\
    mode 5 & box/cavity & 20.703 & 20.698 & 0.026 \\
    mode 6 & box/cavity & 26.088 & 26.081 & 0.027 \\
    \bottomrule
  \end{tabular}
\end{table}

The $17.49$\,GHz mode is a shared artifact of the isotropic sheet term:
the junction patch has one free node whose hat-function gradient carries no
port voltage but stores energy in $S_\Gamma$, with Rayleigh quotient
$1/(L_JC_J)$, or $f_{LC} = 17.60$\,GHz, lowered by the adjacent volume
mass; Palace gives it a port EPR of $2.49\times10^{-8}$. This mechanism is
our inference, supported by a synthetic test; the rank-1 shunt removes the
mode. The fixture is coarse: Palace's own refinements of it move the qubit from
$3.4527$ to $3.9528$\,GHz at $h/2$ ($+14.5\%$) and to $4.1721$\,GHz at
second order on $h$ ($+20.8\%$), and the resonator by $+6.1\%$ at $h/2$, so
the agreement above is a code-to-code comparison on one discretization.

\section{The DeviceLayout Cross-Check}
\label{sec:crosscheck}

\paragraph{The published designs.}
The DeviceLayout.jl blog \citep{devicelayout-blog} reports the initial
design ($c = 620\,\mu$m, $L = 5000\,\mu$m) at $[4.14, 5.591]$\,GHz and
the final design ($c = 1217.678\,\mu$m, $L = 7243.955\,\mu$m) at
$[3.005, 3.982]$\,GHz, reached by PRIMA in nine Palace runs ``within 1\% of
our target frequencies''. The example source \citep{devicelayout-example}
sets the targets to $[3.0, 4.0]$\,GHz, works in the variables $x_1 =
\sqrt{620\,\mu\mathrm{m}/c}$ and $x_2 = 5000\,\mu\mathrm{m}/L$, and stops
when $\mathrm{err} = \frac12\sum_m (f_m - t_m)^2/t_m^2 < 10^{-4}$. Its
frequencies are computed with Palace on DeviceLayout's mesh at mesh order 2
and solver order 2, with a first-order absorbing far field and two
$50\,\Omega$ readout ports.

\paragraph{Reproduction at DeviceLayout's discretization.}
Palace on DeviceLayout's second-order mesh with the example's
configuration file gives $[4.1683, 5.6132]$\,GHz for the initial design
and $[3.0164, 3.9956]$\,GHz for the final one, $0.68\%$ and $0.40\%$
above the published initial pair and $0.38\%$ and $0.34\%$ above the
final pair. The same cells with the PEC-box model used everywhere else
differ from the configuration-file runs by at most $0.003\%$. Neither the
blog nor the example records the Palace version, mesher build or platform
of the published runs (ours are in Section~\ref{sec:repro}). Every published value lies
below ours, as a slightly coarser published mesh would make it, but which
difference causes the $0.34$--$0.68\%$ gap is not diagnosed.

\paragraph{The ladder and its rate.}
For each geometry, Palace runs at order 1 on DeviceLayout's first-order
mesh $h$ and its uniform refinements $h/2$ and $h/4$ (gmsh's eight-way
split, up to $12.6$ million unknowns), and at order 2 on $h$ and $h/2$;
geode-fem runs at order 1 on $h$ and $h/2$, on the same meshes. Every
computed series rises under refinement (Figure~\ref{fig:crosscheck}); this
is an observation, not a guarantee. At the edge of a zero-thickness metal
sheet the field behaves as $E \sim r^{-1/2}$ \citep{meixner1972behavior},
so near the edge the eigenfunctions lie in $H^{s}$ only for $s < 1/2$, and
on uniform meshes the eigenvalue error decays as $h^{2s}$ at any element
order \citep{boffi2010finite}, slower than $O(h)$. The order-1
ladders show pre-asymptotic orders of $0.61$--$0.62$ (qubit) and $0.70$--$0.73$
(resonator).

\paragraph{Extrapolation rule.}
The \emph{estimate} is the mean of two extrapolations at order 1 in the
frequency, $2f_2 - f_1$: one from the order-1 pair ($h/2$, $h/4$) and one
from the order-2 pair ($h$, $h/2$). Order 1 is the fastest rate the theory
allows on uniform meshes, so the rule makes the smallest correction
consistent with it, and its estimates lean low; for the order-2 pair the
order is taken from theory, because two levels give no observed order. The
\emph{band} spans six values, the same two pairs taken three ways (values
for the initial qubit):
\begin{itemize}\setlength{\itemsep}{0pt}
  \item at order 1 in the frequency, the two values of the estimate
    ($4.432$ and $4.455$\,GHz);
  \item at order 1 in the eigenvalue variable $\lambda\propto f^2$,
    $\sqrt{2f_2^2 - f_1^2}$, which lie lower for a series that rises
    under refinement ($4.415$ and $4.451$\,GHz); the order-1 value sets
    the low end, except on family-b resonators, where the order-2 value
    does;
  \item at the observed order of the order-1 ladder's three levels
    ($0.62$ and $0.61$ for the family-a qubits), borrowed by the order-2
    pair on the assumption that both element orders share the
    pre-asymptotic rate, which lie higher ($4.677$ and $4.582$\,GHz); the
    order-1 value sets the high end.
\end{itemize}
The band is an extrapolation range, not a bound. For the initial qubit,
the circuit-model cap of Section~\ref{sec:static} trims it at
$4.477$\,GHz.

The rule is Richardson extrapolation \citep{richardson1911approximate} at
the asymptotic order rather than the observed one. Roache's grid
convergence index \citep{roache1994perspective} scales the correction at
the observed order by a safety factor: on the initial qubit's order-1
ladder, the fine-grid index with a factor of $1.25$ is centred on
$4.155$\,GHz with a half-width of $0.65$\,GHz, which reaches $4.81$\,GHz,
above the circuit-model cap. The rule was fixed after the static route had
been computed; extrapolation at the observed order alone puts the initial
qubit at $4.63$\,GHz, above the cap, and the observed-order values are kept
only as the band's high side. The agreement of the two routes is therefore
a consistency check, not a blind prediction.

\begin{figure}[htbp]
  \centering
  \anvilfig{fig2-crosscheck.pdf}{fig2\_crosscheck.py}
  \caption{Resolution ladders of the two published DeviceLayout.jl
  SingleTransmon designs: Palace at orders 1 and 2 on DeviceLayout's mesh
  $h$ and its uniform refinements, and geode-fem at order 1 on the same
  meshes (crosses); the open diamond, just right of $h$, is Palace at
  DeviceLayout's own discretization, that of the published values
  (dashed). The ``$h\to0$'' slots are categorical: the estimate with its
  band for the gmsh refinement family (a, labelled) and for Palace's
  refinement (b). In (a), the static route of Section~\ref{sec:static}
  (grounded box) and the circuit-model cap, which trims the untrimmed
  bands drawn. Sources:
  \texttt{benchmarks/transmon\_devicelayout\_crosscheck/results.toml} and
  \texttt{benchmarks/transmon\_devicelayout\_controls/results.toml}.}
  \label{fig:crosscheck}
\end{figure}

\begin{table}[htbp]
  \centering
  \caption{Published frequencies, Palace at DeviceLayout's discretization
  (DL), the finest level of Palace's adaptive run (a measured value, still
  rising), and the extrapolated estimates with their bands, in GHz, for
  the two refinement families of DeviceLayout's mesh $h$ (a: gmsh's
  eight-way split; b: Palace's uniform refinement). Percentages are
  relative to the published value. Sources:
  \texttt{benchmarks/transmon\_devicelayout\_crosscheck/results.toml}
  (family a) and \texttt{benchmarks/transmon\_devicelayout\_controls/results.toml}
  (adaptive, family b).}
  \label{tab:crosscheck}
  \footnotesize
  \setlength{\tabcolsep}{3pt}
  \begin{tabular}{@{}lccc@{\hspace{6pt}}ll@{}}
    \toprule
    & & & Adaptive, & \multicolumn{2}{c}{Extrapolated estimate [band]} \\
    \cmidrule(l){5-6}
    & Published & DL & finest level & Family a & Family b \\
    \midrule
    Initial qubit & 4.14 & 4.1683 & 4.4297 ($+7.0\%$) & 4.444 ($+7.3\%$) $[4.415, 4.677]^{\ast}$ & 4.472 ($+8.0\%$) $[4.458, 4.510]$ \\
    Initial resonator & 5.591 & 5.6132 & 5.6986 ($+1.9\%$) & 5.740 ($+2.7\%$) $[5.723, 5.817]$ & 5.776 ($+3.3\%$) $[5.768, 5.800]$ \\
    Final qubit & 3.005 & 3.0164 & 3.2468 ($+8.0\%$) & 3.265 ($+8.7\%$) $[3.236, 3.471]$ & 3.292 ($+9.6\%$) $[3.278, 3.331]$ \\
    Final resonator & 3.982 & 3.9956 & 4.0552 ($+1.8\%$) & 4.088 ($+2.7\%$) $[4.077, 4.139]$ & 4.113 ($+3.3\%$) $[4.106, 4.125]$ \\
    \bottomrule
  \end{tabular}

  \smallskip
  {\footnotesize $^{\ast}$Under the circuit model of
  Section~\ref{sec:static}, the cap of $4.477$\,GHz trims the band to
  $[4.415, 4.477]$.}
\end{table}

\paragraph{Result.}
The initial qubit extrapolates to $4.444$\,GHz and the final qubit to
$3.265$\,GHz, $7.3\%$ and $8.7\%$ above the published values; both
resonators extrapolate to $2.7\%$ above theirs
(Table~\ref{tab:crosscheck}). Every published value lies below its band.
Measured against the estimates of the two families, DeviceLayout's
discretization sits $6.2$--$6.8\%$ and $7.6$--$8.4\%$ low on the two qubits
and $2.2$--$2.8\%$ low on the resonators. With no extrapolation, the
highest uniform levels, order 2 on $h/2$ ($7.0$ million unknowns on the
initial geometry; order 1 on $h/4$ has more unknowns and lower values),
already lie $4.2\%$ and $4.7\%$ above the published qubit values. The
estimates rest on Palace's order-1 $h/4$ and order-2 levels, which
geode-fem did not run. On the levels it did run, geode-fem and Palace
agree on the qubit to $10^{-4}\,\%$ on both geometries. On the resonator
they differ by $0.03\%$ at $h$, where Palace adds $4{,}946$ (initial) and
$5{,}944$ (final) elements by local bisection, and not at all in the six
archived decimals at $h/2$, where Palace's logs report no bisection and
the two codes differ only in the junction term
(Section~\ref{sec:junction}; the counts are in the logs). Three checks
test the estimates, and all three concern the qubit.

\paragraph{A second refinement family.}
Family~b refines the same $h$ mesh with Palace's own uniform refinement
instead of gmsh's split (family~a). It first bisects where a sheet is one
element wide (Section~\ref{sec:agreement}), so it is graded slightly
toward the sheets, a plausible reason for its higher observed orders,
$0.91$ (initial) and $0.90$ (final) against family~a's $0.62$ and $0.61$.
Its $h/2$ qubit is $3.9469$\,GHz against gmsh's $3.8771$\,GHz, and on the
fixture the two refinements give $3.9528$ and $3.8939$\,GHz, so the
refinement scheme accounts for $1.5$--$1.8\%$ of the $h/2$ level and the
choice of $h$ mesh for only $0.15$--$0.43\%$. Its qubit estimates lie
$0.64\%$ and $0.83\%$ above family~a's, and every family-b band lies
inside the family-a band (Table~\ref{tab:crosscheck}). We therefore quote
the initial qubit as $4.444$--$4.472$\,GHz, $7.3$--$8.0\%$ above the
published value, and the final qubit as $3.265$--$3.292$\,GHz, $+8.7\%$ to
$+9.6\%$ over the published $3.005$\,GHz.

\paragraph{Adaptive refinement.}
A further check comes from Palace's adaptive mesh refinement, run on each
geometry from the same $h$ mesh at order 2, with the error indicator summed
over the qubit and resonator modes, D\"orfler marking at $0.7$ and
nonconformal refinement (Figure~\ref{fig:initialq} for the initial
geometry). On the initial
geometry, in five solves from $0.95$ to $11.0$ million unknowns, the qubit
rises from $4.1683$ to $4.4297$\,GHz, $0.32\%$ below the estimate and inside
the band; the third solve, at $2.2$ million unknowns, already exceeds the
highest uniform level ($4.3118$\,GHz at $7.0$ million). On the final
geometry, from $1.17$ to $14.6$ million unknowns, the qubit rises from
$3.0164$ to $3.2468$\,GHz, $8.0\%$ above the published $3.005$\,GHz, and
the third solve again exceeds the highest uniform level ($3.1884$ against
$3.1464$\,GHz). The finest final level lies inside the band, $0.56\%$ and
$1.39\%$ below the two family estimates. Both finest levels are measured
values, not extrapolations, and neither run reached a plateau: each stopped
when its last solve passed the run's size limit of $10^7$ unknowns, not on
its tolerance, and the last step still moved the qubit by $+0.55\%$ (initial)
and $+0.63\%$ (final). A three-point fit $f_\infty - CN^{-\alpha}$ through
the last three solves gives $4.4552$\,GHz ($\alpha = 0.80$) and
$3.2672$\,GHz ($\alpha = 0.83$), but $+0.1\%$ on the finest level moves
them by $24$ and $16$\,MHz, so they are ill-conditioned estimates. The
adaptive resonators rise more slowly and still lie below their bands, so
the adaptive runs check the qubit estimates only.

\paragraph{A common-limit fit.}
The last check fits the
order-1 ($h/2$, $h/4$) and order-2 ($h$, $h/2$) series of each ladder with
one common limit and order. This fit is exactly determined; it lies
$0.5$--$1.3\%$ above the rule's estimate on every first-family series, as
expected of a rule that leans low, and $+0.1\%$ on the order-2 $h/2$ level
alone moves the initial-qubit fit by $20$\,MHz.

\begin{figure}[htbp]
  \centering
  \anvilfig{fig3-initial-qubit.pdf}{fig3\_initial\_qubit.py}
  \caption{Checks of the extrapolated estimate on the initial geometry.
  Left: every Palace level against the number of N\'ed\'elec unknowns, for
  the two refinement families of $h$ (filled: gmsh's split, as in
  Figure~\ref{fig:crosscheck}; open: Palace's) and Palace adaptive
  refinement from $h$ at order 2 (diamonds). Right: the estimates of the
  two families with their bands and, for the qubit only, the adaptive
  run's ill-conditioned fit (the resonator fit is unreliable). In (a), the
  circuit-model cap. Sources:
  \texttt{benchmarks/transmon\_devicelayout\_crosscheck/results.toml} and
  \texttt{benchmarks/transmon\_devicelayout\_controls/results.toml}.}
  \label{fig:initialq}
\end{figure}

\section{A Static Route: A Bounded Island Capacitance}
\label{sec:static}

\paragraph{The bound.}
For the reduced electrostatic problem (island at potential 1, ground and
resonator at 0, feedline floating with zero net charge, and zero normal
flux on the vacuum box), the potential
$\phi$ satisfies $0 \le \phi \le 1$ by the maximum principle, and
$\Cisl = \int \varepsilon\,\nabla\phi\cdot\nabla\phi$ with $\varepsilon
= \varepsilon_0\varepsilon_r$. Every conforming finite-element energy is an
upper bound on $\Cisl$, because $\Cisl$ is the minimum of the energy
over admissible potentials. For the lower bound, take any piecewise
polynomial field $D$ on the mesh cut along the zero-thickness sheets.
Integrating by parts and using $0\le\phi\le1$ on every term gives
\begin{equation}
  \int D\cdot\nabla\phi \;\ge\; N(D) \;=\; Q_{\mathrm{island}}(D)
  - \lvert Q_{\mathrm{feed}}(D)\rvert
  - \lVert \operatorname{div} D\rVert_{L^1}
  - \lVert [D\cdot n]\rVert_{L^1(\mathcal F)}
  - \lVert D\cdot n\rVert_{L^1(\mathcal B)},
  \label{eq:cert-n}
\end{equation}
where $Q$ are the fluxes of $D$ into the two conductors (for the exact
field $D = \varepsilon\nabla\phi$, the negative of the physical
displacement, they are the conductor charges),
$[D\cdot n]$ is the normal jump on the interior faces $\mathcal F$ off
the sheets, and $D\cdot n$ on the box faces $\mathcal B$ is the one-sided
flux, which the zero-flux condition requires of the exact field but not
of $D$. Cauchy--Schwarz, $\bigl(\int D\cdot\nabla\phi\bigr)^2 \le
\int D\cdot\varepsilon^{-1}D \int\varepsilon\nabla\phi\cdot\nabla\phi$,
then gives, whenever $N(D) > 0$,
\begin{equation}
  \Cisl \;\ge\; \frac{N(D)^2}{\int D\cdot\varepsilon^{-1}D}.
  \label{eq:cert}
\end{equation}
The bound holds for any $D$, so solver tolerance can only weaken it. Each
$L^1$ norm is bounded above element by element and face by face,
$\lVert p\rVert_{L^1(K)} \le |K|^{1/2}\lVert p\rVert_{L^2(K)}$, with
$\lVert p\rVert_{L^2}^2$ integrated exactly by quadrature. We
take $D$ from the P2 potential, equilibrated patch by patch into
Raviart--Thomas fluxes of order one to three, with an independent
Crouzeix--Raviart flux as a check. The two bounds hold by theorem, for the
exact arithmetic of the discretization on the polygonal geometry of the
mesh, with lower bounds floored and upper bounds ceiled; the implementation
is checked by tests, not formally verified. The tests compare it with an
exact parallel plate and an analytic coaxial bracket, and mutation tests
reject a grounded-feedline or jump-free lower bound. The bound is evaluated
in plain f64, not interval arithmetic; the worst-case summation error is about
$10^{-10}$ relative, far below the rounding. Refinement does not move the
polygonal surfaces, so all meshes of one geometry bound the same exact
capacitance, and a family bracket is the largest lower bound and the
smallest upper bound.

\paragraph{Two box models.}
The bound above leaves the vacuum box natural; the eigenproblem makes it
PEC (Section~\ref{sec:geometry}). Constraining the potential on the box,
to zero or to a floating constant, shrinks the set of admissible
potentials, so it can only raise $\Cisl$: the natural-box lower bound, and
the cap below that uses only it, hold for the PEC box too. For an upper
bound we re-solve the conforming P2 problems on the same meshes with the
box grounded.

\paragraph{The unmodified device.}
Uniform refinement of the fixture with P1 and P2 elements gives the upper
bounds of Figure~\ref{fig:convergence}: with the scalar model, from
$137.706752$\,fF at P1 on $h$ to $83.015285$\,fF at P2 on $h/2$, and with
the tensor from $136.684731$ to $85.008425$\,fF. On adaptively refined meshes the
lower bound closes the gap: the island capacitance of the unmodified
initial design is bounded by $[77.40, 77.61]$\,fF with the scalar
$\bar\varepsilon$ and by $[79.56, 79.76]$\,fF with the sapphire tensor
(Figure~\ref{fig:convergence}). The bounded tensor-minus-scalar
difference, computed from the unrounded bounds, is $1.95$--$2.35$\,fF, or
$2.5$--$3.1\%$ of the scalar value. With the box grounded, the tensor
$\Cisl$ is bounded by $[79.56, 79.77]$\,fF. On the finest mesh, grounding
the box raises the conforming value by $0.0065$\,fF, $8\times10^{-5}$
relative; that is an estimate of the effect of the box model, whose size is
bounded only by $[0, 0.21]$\,fF; the same bracket covers a floating box,
which lies between the two. With $C_J$ added, the transmon's
$C_\Sigma = \Cisl + C_J$ is bounded by $[85.06, 85.27]$\,fF for the
eigenproblem's grounded box (tensor; $[85.06, 85.26]$\,fF natural). The
fixture itself overstates the tensor $\Cisl$ by at least $56.92$\,fF, which is at least $71.3\%$ of
the exact value, and the scalar one by at least $77.4\%$.

\begin{figure}[htbp]
  \centering
  \anvilfig{fig4-convergence.pdf}{fig4\_convergence.py}
  \caption{Mesh convergence of the island capacitance $\Cisl$ of the
  unmodified device under uniform refinement of the fixture, with the scalar trace-averaged
  permittivity (blue) and the sapphire tensor (orange). Every marker is a
  conforming upper bound on the exact capacitance of the meshed polygonal
  geometry. The short bars are the two-sided bounds from
  adapted meshes, each drawn at the tetrahedron count of the mesh that
  sets its lower bound and magnified in the inset; the inset title gives the bounded
  tensor-minus-scalar gap. All values are for the natural box. Sources: \texttt{c\_sigma\_convergence.toml},
  \texttt{c\_sigma\_bracket.toml} and
  \texttt{c\_sigma\_bracket\_tensor.toml}.}
  \label{fig:convergence}
\end{figure}

\paragraph{From capacitance to the qubit frequency.}
Reading the grounded-box tensor bound through the two-node circuit of
Section~\ref{sec:junction}, over every corner of the $\Cisl$, $L_s$ and
$k$ bands, with the resonator at its extrapolated estimate, gives the
initial dressed qubit in $[4.449, 4.458]$\,GHz; the natural-box bound moves
its low end by $0.3$\,MHz. That range lies inside the eigenmode band of
Section~\ref{sec:crosscheck} and inside the range of the two family
estimates, and the first family's estimate is $0.22\%$ below its
midpoint, a consistency check for the reason given in
Section~\ref{sec:crosscheck}. The range is an estimate, because $L_s$, $k$
and the resonator estimate are.

\paragraph{A model-dependent cap.}
Under the circuit model, a series inductance and a coupling to a higher
mode can only lower the qubit, so the lower bound $\Cisl \ge 79.56$\,fF,
that is $C_\Sigma \ge 85.06$\,fF, caps it at $1/(2\pi\sqrt{L_J C_\Sigma})
= 4.477$\,GHz. The cap is rigorous only given that model (one island node
with $L_J \parallel C_J$, $L_s \ge 0$, coupling only to modes above the
qubit, a quasi-static $C_\Sigma$) and rests on the fixture, not on the
ladder meshes, which agree with it to $0.02\%$ at $h$, so it is
model-dependent and not in the Bounded tier. It does not depend on the
extrapolation rule. Every computed level and both family estimates lie
below it; the first family's common-limit fit, $4.4802$\,GHz, exceeds it by
$3.6$\,MHz, within that fit's conditioning. With no extrapolation, the
adaptive level and the cap put the initial qubit between $4.4297$ and
$4.477$\,GHz, $+7.0\%$ to $+8.1\%$ over the published value; the lower end
is a measured level, the upper end holds only under the circuit model.

\paragraph{Why the qubit is approached from below.}
The qubit mode is quasi-static, and the quasi-static part of a first-order
N\'ed\'elec field is the gradient of a P1 potential, which can only
overstate the capacitance, so the computed qubit should approach its limit
from below: on the fixture the P1 bound ($136.68$\,fF, with $C_J$ added)
overstates $C_\Sigma$ by a factor of $1.67$, and the family estimates exceed
the fixture's order-1 qubit by $1.66$--$1.68$ in $f^2$. This is a heuristic, not a proof: it omits $L_s$,
the coupling, the rotational field and the distributed resonator.

\section{Re-Optimization and Its Controls}
\label{sec:reopt}

At the extrapolated frequencies the published final design misses its
targets. We repeat DeviceLayout's optimization against its own
discretization, to compare optimizers in one environment, and against the
extrapolated frequencies, to find the design that meets the targets.

\paragraph{The capacitor\_length derivative.}
DeviceLayout builds the island upward from the fixed junction, so a change
of $c$ translates every node above the island and stretches the island
band. geode-fem turns this motion into a node velocity
$\partial X/\partial c$, checks that every metal outline node in a
stretched band moves along its own edge ($3{,}146$ nodes on the initial
mesh), and evaluates Eq.~\eqref{eq:hf-eig} with the sapphire tensor for both modes.
The contraction takes $0.44$\,s. On the initial mesh it matches a central
difference of two re-solves of the moved mesh ($\pm0.5\,\mu$m) to every
printed digit: $-2.443227\times10^{-3}$\,GHz/$\mu$m for the qubit and
$-3.964\times10^{-6}$\,GHz/$\mu$m for the resonator.

\paragraph{The total\_length column.}
DeviceLayout re-lays the meander when $L$ changes, so no fixed-topology
node motion represents it. A transmission-line model holds $f_r(L+\delta)$
constant and takes $\partial f_q/\partial L = 0$, with $\delta =
216\,\mu$m from one finite difference of two geode-fem order-1 solves
($L = 5000$ and $5250\,\mu$m). $\delta$ is weakly determined: remesh noise
of about $0.15\%$ in either $f_r$ moves it by about $\pm175\,\mu$m. The
model enters only the Jacobian, and every iterate is evaluated on the true
objective, so a poorer Jacobian costs steps, not accuracy.

\paragraph{Newton on the high-fidelity objective.}
DeviceLayout re-meshes every design, so the first-order objective is not
smooth in the parameters: from $c = 620$ to $650\,\mu$m the remeshed
first-order difference of the qubit is $-5.23\times10^{-3}$\,GHz/$\mu$m
against a fixed-topology derivative of $-2.28\times10^{-3}$ averaged over
the interval (the mean of its end values), because remeshing alone moves
the qubit by about $2.6\%$, more than the $2.0\%$ the step moves it. The
iteration therefore runs Newton directly on the high-fidelity objective
$F$, in DeviceLayout's variables $(x_1, x_2)$, in which the frequencies are
close to linear and nearly decoupled: a probe in either variable moves its
own mode by $4.5\%$ and the other by at most $0.05\%$. The $c$ column of
the Jacobian is geode-fem's first-order derivative scaled by the fidelity
ratio $S = F/f_h$ at the same design. Each step costs one geode-fem
first-order solve with its gradient ($19$--$25$\,s) and one high-fidelity
evaluation; the geode-fem solves of the one-time $\delta$ calibration are
outside the Palace-run count.

\begin{table}[htbp]
  \centering
  \caption{Palace runs to DeviceLayout's criterion ($\mathrm{err} <
  10^{-4}$) at its discretization, all in our Palace environment except the
  published PRIMA run. Every Palace run a method uses counts, with the
  initial design as run~1 and probes included, among them the probe that
  sets the Palace $\delta$ (Broyden's second probe, shared by the two
  rows); geode-fem solves are not counted. Each count is one trajectory.
  PRIMA runs at \texttt{rhobeg} $= 0.2$ unless stated; the Newton variants
  use the transmission-line $L$ column.
  Sources: the \texttt{optimize} and \texttt{controls} artifacts under
  \texttt{benchmarks/transmon\_devicelayout\_*/}.}
  \label{tab:controls}
  \footnotesize
  \begin{tabular}{@{}lllccc@{}}
    \toprule
    Method & $c$ column & $\delta$ from & Runs & err after first step & Design at criterion ($\mu$m) \\
    \midrule
    PRIMA, published & none & none & 9 & unpublished & (1217.678, 7243.955) \\
    PRIMA, our environment & none & none & 10 & n/a & (1209.125, 7207.551) \\
    \quad \texttt{rhobeg} $0.1$ & none & none & 10 & n/a & (1235.705, 7216.766) \\
    \quad \texttt{rhobeg} $(x_u - x_l)/2$ & none & none & 12 & n/a & (1245.900, 7291.401) \\
    Newton & shape derivative & geode-fem & 3 & $7.9\times10^{-4}$ & (1225.346, 7223.250) \\
    Newton & scaling law & geode-fem & 3 & $1.7\times10^{-4}$ & (1232.664, 7221.270) \\
    Newton & scaling law & Palace & 4 & $2.1\times10^{-4}$ & (1229.026, 7225.712) \\
    Broyden & two probes & none & 5 & $6.2\times10^{-4}$ & (1224.847, 7232.447) \\
    \bottomrule
  \end{tabular}
\end{table}

\begin{figure}[!tbp]
  \centering
  \anvilfig{fig5-reoptimization.pdf}{fig5\_reoptimization.py}
  \caption{The re-optimization and its controls. (a) The design plane near
  the optimum: the published final design, PRIMA's end point in our
  environment, the Newton iterates against DeviceLayout's discretization
  (numbered: shape derivative; the first legs start off the plot), and the
  steps against the extrapolated frequencies, with the step-2 qubit band
  mapped linearly to capacitor\_length (an estimate). (b) DeviceLayout's
  err against Palace evaluations at its discretization in our environment,
  probes counted (the Palace-$\delta$ Newton's run 2 is the $L$ probe it
  shares with Broyden), PRIMA at \texttt{rhobeg} $= 0.2$; the published PRIMA
  point is at its 9 runs and the err of its published frequencies,
  $1.151\times10^{-5}$ (DeviceLayout's environment). This counts Palace
  evaluations, not total cost: each shape-derivative Newton point also
  costs one geode-fem solve. (c, d) Extrapolated
  estimates and bands against the targets, with step~2 in both refinement
  families. Sources: the \texttt{optimize}, \texttt{controls} and
  \texttt{crosscheck} artifacts under
  \texttt{benchmarks/transmon\_devicelayout\_*/}.}
  \label{fig:reopt}
\end{figure}

\paragraph{Against DeviceLayout's discretization.}
Table~\ref{tab:controls} and Figure~\ref{fig:reopt}(b) compare the
optimizers on one objective. Newton
with the shape-derivative column meets the criterion at its second step,
after 3 Palace runs, and its three geode-fem solves add about $70$\,s,
against about $4$\,min per Palace run on 32 ranks; a third step reaches
$\mathrm{err} = 1.4\times10^{-7}$ at $(1236.360, 7224.900)\,\mu$m. Two
gradient-free controls test where the 3 runs come from. The first replaces
the $c$ column by the scaling law $\partial F_q/\partial x_1 = F_q/x_1$,
which follows from $f_q \propto 1/\sqrt{C_\Sigma}$ with $C_\Sigma$ taken
proportional to $c$. With the
geode-fem $\delta$ it also takes 3 runs, and its first step has an err
$4.7$ times lower (Table~\ref{tab:controls}), about $2.2$ times closer in
frequency. With $\delta = 563.879\,\mu$m taken
instead from a Palace probe ($c$ fixed, $L$ from $5000$ to
$5263\,\mu$m), so that no geode-fem input remains, it takes 4 runs: the
same two steps, plus the probe. This $\delta$ is $2.6$ times the geode-fem
one, more than the $\pm175\,\mu$m noise estimate, and costs no extra step.
The second control, Broyden's method \citep{broyden1965class} from two
forward finite-difference probes, takes 5.

\paragraph{PRIMA in our environment.}
DeviceLayout's own PRIMA call,
re-run verbatim in our environment, takes 10: runs 1--5 are its initial
interpolation set of $2n+1 = 5$ points \citep{powell2009bobyqa}, and run 9
misses the criterion by $15\%$ ($\mathrm{err} = 1.150\times10^{-4}$). In
this environment the published final design already meets the criterion
($\mathrm{err} = 1.5\times10^{-5}$), so the environment moves PRIMA's
trajectory, not the attainability of the target; the published count of 9
is from DeviceLayout's environment. With the initial trust radius
\texttt{rhobeg} at $0.1$ instead of $0.2$, PRIMA also takes 10 runs, and
at the largest radius BOBYQA accepts on DeviceLayout's box, $(x_u -
x_l)/2$ (one ulp below $0.4$, which PRIMA rejects), it takes 12. Each setting is a single run. In one environment, then, the
comparison is 3 Newton runs against 10 to 12, or 4 with no geode-fem input.

\paragraph{Where the first steps land.}
The shape-derivative and Broyden first steps land at $c = 1338.338$ and
$1320.522\,\mu$m, past about $1230\,\mu$m at the criterion, because both
follow a local slope, $\partial F_q/\partial x_1 = 3.66$ and $3.71$\,GHz,
within about $1.5\%$ of each other and $11$--$12\%$ shallower than the
scaling law's $F_q/x_1 = 4.168$\,GHz.

\paragraph{Against the extrapolated frequencies.}
Starting from the published final design, each Newton step uses the
extrapolated estimate and the geode-fem derivative at the current design,
and each new design is evaluated by a full ladder with the rule of
Section~\ref{sec:crosscheck} (Table~\ref{tab:newton-conv}).
The first overshoots the qubit by $-1.09\%$; the second lands at $+0.16\%$
and $-0.04\%$ of the targets, with both targets inside the bands. On both
ladders geode-fem and Palace agree on the qubit to $2\times10^{-4}\,\%$.
One scaling-law step from the same start, evaluated by the Palace half of
the ladder, lands $3.6$ and $8.2\,\mu$m short of step~2 and at $+0.13\%$
and $+0.12\%$ of the targets: one gradient-free step does what the two
shape-derivative steps did.

\begin{table}[htbp]
  \centering
  \caption{Steps against the extrapolated frequencies, in GHz. Each design
  is evaluated by a full resolution ladder (the scaling-law step and the
  second family by Palace alone) and the rule of
  Section~\ref{sec:crosscheck}; values are estimates with their bands.
  Family b refines step 2's $h$ mesh with Palace's uniform refinement
  instead of gmsh's. Sources: the \texttt{optimize}, \texttt{controls} and
  (published final) \texttt{crosscheck} artifacts under
  \texttt{benchmarks/transmon\_devicelayout\_*/}.}
  \label{tab:newton-conv}
  \footnotesize
  \begin{tabular}{@{}lcccc@{}}
    \toprule
    Design & $c$ ($\mu$m) & $L$ ($\mu$m) & $f_q$ [band] & $f_r$ [band] \\
    \midrule
    Published final & 1217.678 & 7243.955 & 3.265 $[3.236, 3.471]$ & 4.088 $[4.077, 4.139]$ \\
    Step 1 & 1481.657 & 7407.945 & 2.967 $[2.938, 3.188]$ & 4.004 $[3.993, 4.055]$ \\
    Step 2 & 1446.078 & 7416.352 & 3.005 $[2.979, 3.221]$ & 3.998 $[3.987, 4.067]$ \\
    Step 2, family b & 1446.078 & 7416.352 & 3.020 $[3.003, 3.101]$ & 4.022 $[4.016, 4.043]$ \\
    Scaling-law step & 1442.479 & 7408.181 & 3.004 $[2.977, 3.235]$ & 4.005 $[3.994, 4.055]$ \\
    \bottomrule
  \end{tabular}
\end{table}

\paragraph{The re-optimized design.}
Step 2 is the design whose extrapolated frequencies sit on the targets:
$c = 1446.078\,\mu$m, $18.8\%$ longer than published, with a resonator
$2.4\%$ longer ($+172\,\mu$m). The second refinement family at step 2 gives
$3.0203$ $[3.0033, 3.1014]$\,GHz for the qubit and $4.0225$\,GHz for the
resonator, $0.52\%$ and $0.60\%$ above the first family, the same spread as
on the initial geometry; both targets lie just below its band, so on family
b step 2 is slightly short. Mapped to $3.0$\,GHz by the scaling law $f_q
\propto c^{-1/2}$, the two families' qubit estimates call for $c = 1451$ and
$1466\,\mu$m, $+19.1\%$ and $+20.4\%$ over the published length. Their
bands map to $+17\%$ to $+37\%$ (family a) and $+19\%$ to $+27\%$
(family b) under the scaling law, and to $+17\%$ to $+39\%$ ($c$ from about
$1423$ to $1690\,\mu$m) and $+19\%$ to $+28\%$ linearly through $dF/dc
\approx S\,\partial f_h/\partial c = -9.05\times10^{-4}$\,GHz/$\mu$m at
the step-1 design, centred on the target. Under either mapping the first
range contains the second. Since the estimates lean low, the common-limit
fit, which puts step~2 at $3.039$ and $4.019$\,GHz, would call for a longer
capacitor still.

\section{How geode-fem Was Designed and Built}
\label{sec:built}

geode-fem implements the whiteroom operator specification of Crutcher
Dunnavant \citep{dunnavant2026whiteroom}, whose four layers, L1--L4,
restate Palace's cited source as pure functions, unfold those into
compositions of base algebraic primitives, recast them where possible as
global tensor operations, and express them in a small graph-evaluation
calculus. He also gave design
input on the solver's operator layer and contributed code: the migration to
Rust 2024, the Metal (Apple GPU) backend, the present layout of the core
crate and the examples, and the geode-util crate. The rest of the code, the
experiments and the text of this paper were produced by AI models
working as role-separated agents: Anthropic's Claude (Opus, Fable and
Sonnet models, through Claude Code) and OpenAI's GPT Sol and GPT Astra
models (via Codex). The agents were coordinated by Loom \citep{loom}, an
issue, worktree and pull-request workflow, and by Squad, a message room
the agents share.
Nearly every change was reviewed by a separate agent session. The agent
reviews of this paper recomputed its numbers from the archived artifacts,
reproduced runs, and mutation-tested the scripts that check the text and
the figures against those artifacts. The human author directed the
research, made every scoping and claim decision, and takes responsibility
for the content.

\section{Discussion}
\label{sec:discussion}

\paragraph{What is established, and how.}
Table~\ref{tab:tiers} sorts the results by the kind of evidence behind
them. The brackets are the only bounds, and the cap follows from a bracket
only under a circuit model.

\begin{table}[htbp]
  \centering
  \caption{Evidence tiers. Bounded: a two-sided bound that holds by theorem
  for the meshed polygonal geometry. Model-dependent: a rigorous consequence of a
  bound under a stated circuit model. Measured: a solve on a stated
  discretization. Estimated: an extrapolation with a stated rule.
  Consistency: an agreement that is not independent evidence.}
  \label{tab:tiers}
  \footnotesize
  \begin{tabular}{@{}>{\raggedright\arraybackslash}p{0.17\linewidth}>{\raggedright\arraybackslash}p{0.77\linewidth}@{}}
    \toprule
    Tier & Results \\
    \midrule
    Bounded & Island capacitance $\Cisl$ of the unmodified device, without $C_J$, natural box: $[77.40, 77.61]$\,fF (scalar), $[79.56, 79.76]$\,fF (tensor); grounded box, the eigenproblem's model: $[79.56, 79.77]$\,fF (tensor), so $C_\Sigma = \Cisl + C_J$ in $[85.06, 85.27]$\,fF and $e^2/2hC_\Sigma$ in $[0.2271, 0.2278]$\,GHz; grounded minus natural in $[0, 0.21]$\,fF; tensor-minus-scalar gap $2.5$--$3.1\%$ \\
    Model-dependent & The initial-qubit cap of $4.477$\,GHz $= 1/(2\pi\sqrt{L_J C_\Sigma})$ at the lower bound on $C_\Sigma$: two-node circuit, quasi-static $C_\Sigma$, bounded on the fixture's polygonal geometry, not on the ladder meshes; with the adaptive level, $[4.4297, 4.477]$\,GHz; the same $e^2/2hC_\Sigma$ read as the transmon's charging energy $E_C/h$ \\
    Measured & Same-mesh agreement (Table~\ref{tab:agreement}); the reproduction within $0.7\%$; every ladder and adaptive level (adaptive qubits $+7.0\%$ and $+8.0\%$, still rising); the gradient check; Palace runs to DeviceLayout's criterion, one trajectory each (Table~\ref{tab:controls}) \\
    Estimated & Extrapolated frequencies and bands (Table~\ref{tab:crosscheck}); the adaptive and common-limit fits; the static-route qubit range; $L_s$ and $k$; the box-model effect on $\Cisl$, $8\times10^{-5}$; the re-optimized capacitor, $+19$ to $+20\%$ over two families and about $+17$ to $+39\%$ within the band \\
    Consistency & The two routes agree to $0.22\%$ on the initial qubit; the extrapolation rule was fixed after the static route was known, so this is not a blind prediction \\
    Not claimed & Agreement with measurement; a verified hit of $3.0$/$4.0$\,GHz in the $h \to 0$ limit; a global optimum; an advantage of the shape derivative on this problem; eigenmode-participation (EPR) derivatives, which need eigenvector derivatives \citep{nelson1976simplified}; a speed advantage over Palace \\
    \bottomrule
  \end{tabular}
\end{table}

\paragraph{Why the shape derivative gives no advantage here.}
The derivative is used at every step of the re-optimization, but three
properties of this benchmark remove its advantage. First, the problem has
two nearly decoupled parameters, close to linear in DeviceLayout's
variables, so Newton's method converges with any reasonable Jacobian: the
shape derivative, a two-probe secant, or the scaling law $f_q \propto x_1$,
which costs nothing. Second, a finite-difference gradient costs only $N+1 =
3$ solves. Third, the coarse-mesh, fixed-topology derivative is the slope
of a different discrete problem from the one the iteration evaluates
(DeviceLayout remeshes every design and Palace evaluates it at second
order); once scaled by the fidelity ratio $S$ it gives about what
Broyden's probe gives (Section~\ref{sec:reopt}). That the derivative is exact for its own
discrete problem and cheap ($0.44$\,s) does not change the number of
Palace runs.

Derivatives pay off in other regimes. With $N$ parameters, a discrete
adjoint gives all $N$ derivatives of one objective for about one extra
solve, where finite differences need $N$ more; coupled targets with no
per-parameter scaling law need a computed Jacobian; and freeform shapes
have many parameters by construction. A companion paper on adjoint shape
optimization of a conformal patch antenna \citep{walters2026conformal}
measures this: at 38 shape parameters the adjoint needs at least $11$
times fewer solves than the best gradient-free run, tuned settings
included, to reach its
$-9.5$\,dB target, the ratio falls to $1.1$--$2.4$ at 4 parameters, and
a tuned full-covariance CMA-ES comes within a factor of $7.7$ at $-9$\,dB.

\paragraph{Threats and limitations.}
Agreement between two codes is evidence against implementation error, not
against shared modelling choices such as the port sheet behind the
$17.49$\,GHz mode. The extrapolations assume the rate the sheet-edge
singularity allows on uniform meshes, their high sides are wide, neither
adaptive run reached a plateau, and $L_s$ and $k$ were extrapolated from
two or three discretizations. Every limit $h \to 0$ here is a limit of
DeviceLayout's model (zero-thickness PEC sheets, lossless and linear, open
readout ports in the ladders); film thickness and kinetic inductance remove
the edge singularity and move the frequencies by amounts not measured
here. The element edges
along the sheet outlines are mostly $10$--$12\,\mu$m long on $h$ and so
$2.5$--$3\,\mu$m on $h/4$, more than ten times a $0.2\,\mu$m film: every
mesh used here sees the zero-thickness edge, and the shift measured here
is the discretization error within the model, for one device.

\paragraph{What the shift means for a design.}
Table~\ref{tab:design} reads the initial qubit as a total capacitance. The
bound gives $E_C/h = e^2/2hC_\Sigma$ in $[0.2271, 0.2278]$\,GHz for the
grounded box, Bounded as a function of $C_\Sigma$ and model-dependent as
the transmon's charging energy (a single-island lumped picture). Under the
two-node circuit with the $L_s$ and $k$ bands, the published $4.14$\,GHz
implies $C_\Sigma$ in $[98.59, 98.72]$\,fF, $15.6$--$16.1\%$ above the
bound, so $E_C$ and the anharmonicity, close to $-E_C$
\citep{koch2007charge}, are larger than the published discretization
implies by as much. In the bare LC reading at fixed $L_J$ ($L_s = k = 0$),
which gives a larger $C_\Sigma$ than the two-node one for every frequency,
the published $99.45$\,fF exceeds the bound by $16.6$--$17.0\%$, the
static-route reading by $15.5$--$15.9\%$ and the adaptive-level reading by
$14.5\%$. Of that factor of $1.166$--$1.169$, the measured adaptive level
accounts for $1.145$, the gap between that still-rising level and the
static route for $0.9$--$1.3\%$ more, and $L_s$ and $k$, which the bare
reading omits, for the remaining $0.6$--$1.2\%$. The transition $f_{01}$
rises by a little less than the linear mode, about $6.6\%$ at the adaptive
level. $L_J$ can retune the qubit frequency but not $E_C$, which the
capacitance alone sets.

\begin{table}[htbp]
  \centering
  \caption{The initial qubit read as a total capacitance $C_\Sigma$, with
  $E_C/h = e^2/2hC_\Sigma$; ``bare'' is the LC reading at fixed $L_J$
  ($L_s = k = 0$), ``two-node'' the circuit of
  Section~\ref{sec:junction} with the $L_s$ and $k$ bands, given first for
  the published value. The bound is for
  the grounded box. Sources: \texttt{c\_sigma\_bracket\_grounded\_box.toml},
  \texttt{benchmarks/transmon\_loop\_inductance/results.toml} and the
  adaptive run of Table~\ref{tab:crosscheck}.}
  \label{tab:design}
  \footnotesize
  \begin{tabular}{@{}llccc@{}}
    \toprule
    Source & Tier & $f_q$ (GHz) & $C_\Sigma$ (fF) & $E_C/h$ (GHz) \\
    \midrule
    Bound & Bounded & -- & $[85.06, 85.27]$ & $[0.2271, 0.2278]$ \\
    Static route & Estimated & $[4.449, 4.458]$ & $85.8$--$86.1$ (bare) & $0.225$--$0.226$ (bare) \\
    Adaptive finest level & Measured & $4.4297$ & $86.87$ (bare) & $0.223$ (bare) \\
    Published & -- & $4.14$ & $[98.59, 98.72]$; $99.45$ (bare) & $0.196$; $0.195$ (bare) \\
    \bottomrule
  \end{tabular}
\end{table}

\paragraph{What a practitioner can do.}
The slow rate comes from the sheet edges, and the standard remedy for an
edge singularity is to refine toward it, with a mesh graded toward the
edges or with $hp$ refinement \citep{demkowicz2006computing}, which recover
faster convergence than uniform refinement. Palace's adaptive runs of
Section~\ref{sec:crosscheck} are measured instances, and the gap-refined,
fourth-order setting of \citet{ye2025electromagnetic}, within $0.3\%$ of
their extrapolated values (Section~\ref{sec:related}), is a published one.
A first check is the strip width: on the fixture, $993$ of the $1{,}007$
edges that trigger Palace's bisection lie where a resonator or feedline
strip is one element wide (Section~\ref{sec:agreement}). The same edge fields set the surface participation
that limits coherence, computed with a two-step finite-element simulation
\citep{wang2015surface} or analytical formulas for finite metal thickness
\citep{martinis2022surface}; such loss estimates need the same edge
resolution as the frequencies.

\section{Conclusion}
\label{sec:conclusion}

Palace reproduces the published DeviceLayout.jl SingleTransmon
frequencies at the example's own discretization, which sits $6.2$--$8.4\%$
low on the qubits: adaptive runs put the two qubits $7.0\%$ and $8.0\%$
above their published values with no extrapolation, and the estimates of
two refinement families, which lean low, put them $7.3$--$9.6\%$ above. geode-fem
reproduces Palace's qubit on every level both codes ran, and a two-sided
bracket on the island capacitance, read through a circuit model, is
consistent with that shift and caps the initial qubit under that model. At
the extrapolated frequencies the $3.0$/$4.0$\,GHz design needs a capacitor
about $19$--$20\%$ longer than published, or $17$--$39\%$ within the band.
Newton's method reaches the
example's criterion in 3 Palace runs with geode-fem input (its shape
derivative, or its calibration of the resonator model) and in 4 without,
against 10 to 12 for PRIMA. For design examples that are meant to be
copied, the discretization is part of the result and belongs in the report
beside it.

\section{Reproducibility}
\label{sec:repro}

Every result is read from an archived, scripted artifact.\footnote{The geode-fem source is public at
\url{https://github.com/rjwalters/geode-fem}. Archival DOI to be added at
deposit.}
% TODO(operator): Zenodo DOI
Each caption names its artifact; the circuit terms $L_s$
and $k$ come from \path{benchmarks/transmon_loop_inductance/} and
\path{benchmarks/transmon_coupling/}, the Palace cells are under
\path{reference/fixtures/transmon_palace/}, and the DeviceLayout meshes are
generated by \path{reference/julia/generate_transmon_param.jl}. The versions are geode-fem commit
\texttt{3174015} for the same-mesh agreement, with the commit of each later
result recorded in its artifact, Palace \texttt{fba6a5b}, DeviceLayout.jl
v1.15.0, PRIMA.jl 0.2.4 and Burn 0.21. Palace ran on AWS
Graviton4 hosts (32 and 16 vCPU; the final-design adaptive run on 24
ranks); the capacitance brackets ran on a 28-core, 96\,GiB workstation and
their grounded-box re-solves on a Graviton4 host. geode-fem's
direct solves at $h/2$ peaked at $94$ and $136$\,GB, so only Palace ran
the larger cells. Palace's ladder cells use SLEPc shift-invert at tolerance $10^{-8}$,
\texttt{Target} $3.0$ (initial) or $2.0$\,GHz (final), and its default
direct solver, SuperLU (order-1 $h$), or AMS-preconditioned GMRES at relative
tolerance $10^{-9}$, on 16 (order-1 $h$) or 32 MPI ranks; the fixture ran
on 8, and the runs at DeviceLayout's discretization keep the example's
tolerances ($10^{-6}$ and $10^{-7}$, \texttt{Target} $2.5$\,GHz). geode-fem's Lanczos shifts are
$5$ and $4$\,GHz, with relative residuals below $4\times10^{-8}$.
The $h$, $h/2$ and $h/4$ meshes have $133{,}355$, $1{,}066{,}840$ and
$8{,}534{,}720$ tetrahedra (initial) and $163{,}431$, $1{,}307{,}448$ and
$10{,}459{,}584$ (final). The meshes were made
with gmsh 4.15.2 (\texttt{gmsh\_jll}) under Julia 1.10.11, and Palace was
built against the MPICH of the Rocky Linux 9 image in
\path{reference/palace/docker/Dockerfile}.

\paragraph{Acknowledgments.}
We thank Crutcher Dunnavant for the whiteroom specification, for his
design input to the solver's operator layer and for his contributions to
the codebase, and the Palace and DeviceLayout.jl developers for publishing
the SingleTransmon example and its frequencies. The AI-orchestrated
development of geode-fem and of this paper is described in
Section~\ref{sec:built}.

\bibliographystyle{plainnat}
\bibliography{refs}

\end{document}
